% GENERATED FILE. Edit canonical lessons, then run npm run generate:books.
% canonical-source-hash: fnv1a64:ac74201d1e662bce
% canonical-lessons: TE-C437-sanklistamaina, TE-C437-saralamaina, TE-C437-aspastamaina, TE-C437-tivramaina, TE-C437-nirdharincu, TE-R437-review-430-437, TE-R437-review-final-437

\chapter{Hard or Plain}
\label{ch:te-hard-or-plain}
\hlchaptermodality{437}

\begin{chapteropening}
\textbf{By the end of this chapter:} \emph{I can say what is complicated, plain, vague or serious.}

Complete the last lesson of chapter 437: I can say what is complicated, plain, vague or serious.
\end{chapteropening}

\section[\texorpdfstring{\te{సంక్లిష్టమైన} (saṅkliṣṭamaina)}{సంక్లిష్టమైన (saṅkliṣṭamaina)}]{\te{సంక్లిష్టమైన} (saṅkliṣṭamaina) — complicated; complex}

\label{lesson:TE-C437-sanklistamaina}



\begin{quote}
Before the new one: say the Telugu for to attract, then the Telugu for to divide.
\end{quote}

\subsection*{You'll want to know: \te{సంక్లిష్టమైన}}
\textbf{\te{సంక్లిష్టమైన}} — \emph{saṅkliṣṭamaina} — \textquotedblleft{}complicated; complex\textquotedblright{}.

\begin{quote}
\te{ఆ} \te{సమస్య} \te{చాలా} \te{సంక్లిష్టమైనది}.
\end{quote}

\begin{quote}
\emph{ā samasya cālā saṅkliṣṭamainadi.}
\end{quote}

\begin{quote}
That problem is very complicated.
\end{quote}

\te{సంక్లిష్టత}, complexity.

\begin{cousinweb}
\textbf{\te{సంక్లిష్టమైన}} (\emph{saṅkliṣṭamaina}) is built on Sanskrit \emph{saṁkliṣṭa}, \textquotedblleft{}pressed together; beset with difficulties\textquotedblright{} (Monier-Williams), with Telugu -\te{మైన} for an adjective.
\end{cousinweb}

\begin{grammarlens}[title={What you've built}]
One more word for saying what happens if.
\end{grammarlens}

\subsection*{Your turn}
\begin{itemize}
\raggedright
  \item \emph{Say it:} \emph{saṅkliṣṭamaina}
  \item \emph{Say it:} \emph{saṅkliṣṭamaina}, once more
  \item \emph{Say it:} \emph{ā samasya cālā saṅkliṣṭamainadi.} — That problem is very complicated.
  \item \emph{Recall:} say the Telugu for to attract, then the Telugu for to divide, then say \emph{saṅkliṣṭamaina} again
\end{itemize}

\begin{tcolorbox}[breakable,colback=teal!4,colframe=teal!35!black,title={Before you move on}]
What does \te{సంక్లిష్టమైన} mean? (\textquotedblleft{}Complicated; complex\textquotedblright{}.) Say it once more.
\end{tcolorbox}
\section[\texorpdfstring{\te{సరళమైన} (saraḷamaina)}{సరళమైన (saraḷamaina)}]{\te{సరళమైన} (saraḷamaina) — uncomplicated; plain}

\label{lesson:TE-C437-saralamaina}



\begin{quote}
Before the new one: say the Telugu for to divide, then the Telugu for complicated.
\end{quote}

\subsection*{You'll want to know: \te{సరళమైన}}
\textbf{\te{సరళమైన}} — \emph{saraḷamaina} — \textquotedblleft{}uncomplicated; plain\textquotedblright{}.

\begin{quote}
\te{అతను} \te{సరళమైన} \te{భాషలో} \te{వివరించాడు}.
\end{quote}

\begin{quote}
\emph{atanu saraḷamaina bhāṣalō vivariñcāḍu.}
\end{quote}

\begin{quote}
He explained in plain language.
\end{quote}

\te{సరళంగా}, plainly.

\begin{cousinweb}
\textbf{\te{సరళమైన}} (\emph{saraḷamaina}) is built on Sanskrit \emph{sarala}, \textquotedblleft{}straight; upright, artless\textquotedblright{} (Monier-Williams), with Telugu -\te{మైన}.
\end{cousinweb}

\begin{grammarlens}[title={What you've built}]
One more word for saying what happens if.
\end{grammarlens}

\subsection*{Your turn}
\begin{itemize}
\raggedright
  \item \emph{Say it:} \emph{saraḷamaina}
  \item \emph{Say it:} \emph{saraḷamaina}, once more
  \item \emph{Say it:} \emph{atanu saraḷamaina bhāṣalō vivariñcāḍu.} — He explained in plain language.
  \item \emph{Recall:} say the Telugu for to divide, then the Telugu for complicated, then say \emph{saraḷamaina} again
\end{itemize}

\begin{tcolorbox}[breakable,colback=teal!4,colframe=teal!35!black,title={Before you move on}]
What does \te{సరళమైన} mean? (\textquotedblleft{}Uncomplicated; plain\textquotedblright{}.) Say it once more.
\end{tcolorbox}
\section[\texorpdfstring{\te{అస్పష్టమైన} (aspaṣṭamaina)}{అస్పష్టమైన (aspaṣṭamaina)}]{\te{అస్పష్టమైన} (aspaṣṭamaina) — vague; unclear}

\label{lesson:TE-C437-aspastamaina}



\begin{quote}
Before the new one: say the Telugu for complicated, then the Telugu for uncomplicated.
\end{quote}

\subsection*{You'll want to know: \te{అస్పష్టమైన}}
\textbf{\te{అస్పష్టమైన}} — \emph{aspaṣṭamaina} — \textquotedblleft{}vague; unclear\textquotedblright{}.

\begin{quote}
\te{అతను} \te{ఒక} \te{అస్పష్టమైన} \te{సమాధానం} \te{ఇచ్చాడు}.
\end{quote}

\begin{quote}
\emph{atanu oka aspaṣṭamaina samādhānaṁ iccāḍu.}
\end{quote}

\begin{quote}
He gave a vague answer.
\end{quote}

\te{అ}-, not, + \te{స్పష్టమైన}, clear, which you know.

\begin{cousinweb}
\textbf{\te{అస్పష్టమైన}} (\emph{aspaṣṭamaina}) is built on Sanskrit \emph{aspaṣṭa}, \textquotedblleft{}indistinct\textquotedblright{} (Monier-Williams), with Telugu -\te{మైన}.
\end{cousinweb}

\begin{grammarlens}[title={What you've built}]
One more word for saying what happens if.
\end{grammarlens}

\subsection*{Your turn}
\begin{itemize}
\raggedright
  \item \emph{Say it:} \emph{aspaṣṭamaina}
  \item \emph{Say it:} \emph{aspaṣṭamaina}, once more
  \item \emph{Say it:} \emph{atanu oka aspaṣṭamaina samādhānaṁ iccāḍu.} — He gave a vague answer.
  \item \emph{Recall:} say the Telugu for complicated, then the Telugu for uncomplicated, then say \emph{aspaṣṭamaina} again
\end{itemize}

\begin{tcolorbox}[breakable,colback=teal!4,colframe=teal!35!black,title={Before you move on}]
What does \te{అస్పష్టమైన} mean? (\textquotedblleft{}Vague; unclear\textquotedblright{}.) Say it once more.
\end{tcolorbox}
\section[\texorpdfstring{\te{తీవ్రమైన} (tīvramaina)}{తీవ్రమైన (tīvramaina)}]{\te{తీవ్రమైన} (tīvramaina) — serious; severe}

\label{lesson:TE-C437-tivramaina}



\begin{quote}
Before the new one: say the Telugu for uncomplicated, then the Telugu for vague.
\end{quote}

\subsection*{You'll want to know: \te{తీవ్రమైన}}
\textbf{\te{తీవ్రమైన}} — \emph{tīvramaina} — \textquotedblleft{}serious; severe\textquotedblright{}.

\begin{quote}
\te{ఆ} \te{ప్రాంతంలో} \te{తీవ్రమైన} \te{నీటి} \te{సమస్య} \te{ఉంది}.
\end{quote}

\begin{quote}
\emph{ā prāntaṁlō tīvramaina nīṭi samasya undi.}
\end{quote}

\begin{quote}
There is a serious water problem in that area.
\end{quote}

\te{తీవ్రంగా}, seriously, severely.

\begin{cousinweb}
\textbf{\te{తీవ్రమైన}} (\emph{tīvramaina}) is built on Sanskrit \emph{tīvra}, \textquotedblleft{}strong, severe, violent, intense\textquotedblright{} (Monier-Williams), with Telugu -\te{మైన}.
\end{cousinweb}

\begin{grammarlens}[title={What you've built}]
One more word for saying what happens if.
\end{grammarlens}

\subsection*{Your turn}
\begin{itemize}
\raggedright
  \item \emph{Say it:} \emph{tīvramaina}
  \item \emph{Say it:} \emph{tīvramaina}, once more
  \item \emph{Say it:} \emph{ā prāntaṁlō tīvramaina nīṭi samasya undi.} — There is a serious water problem in that area.
  \item \emph{Recall:} say the Telugu for uncomplicated, then the Telugu for vague, then say \emph{tīvramaina} again
\end{itemize}

\begin{tcolorbox}[breakable,colback=teal!4,colframe=teal!35!black,title={Before you move on}]
What does \te{తీవ్రమైన} mean? (\textquotedblleft{}Serious; severe\textquotedblright{}.) Say it once more.
\end{tcolorbox}
\section[\texorpdfstring{\te{నిర్ధారించు} (nirdhāriñcu)}{నిర్ధారించు (nirdhāriñcu)}]{\te{నిర్ధారించు} (nirdhāriñcu) — to confirm; to ascertain}

\label{lesson:TE-C437-nirdharincu}



\begin{quote}
Before the new one: say the Telugu for vague, then the Telugu for serious.
\end{quote}

\subsection*{You'll want to know: \te{నిర్ధారించు}}
\textbf{\te{నిర్ధారించు}} — \emph{nirdhāriñcu} — \textquotedblleft{}to confirm; to ascertain\textquotedblright{}.

\begin{quote}
\te{పోలీసులు} \te{ఆ} \te{వార్తను} \te{నిర్ధారించారు}.
\end{quote}

\begin{quote}
\emph{pōlīsulu ā vārtanu nirdhāriñcāru.}
\end{quote}

\begin{quote}
The police confirmed that news.
\end{quote}

Past (I …): \te{నిర్ధారించాను}, \emph{nirdhāriñcānu}.

The noun is \te{నిర్ధారణ}, confirmation.

\begin{cousinweb}
\textbf{\te{నిర్ధారించు}} (\emph{nirdhāriñcu}) is built on Sanskrit \emph{nirdhāraṇa}, \textquotedblleft{}specifying, defining, settling, certainty\textquotedblright{} (Monier-Williams), with the Telugu verb ending -\te{ఇంచు}.
\end{cousinweb}

\begin{grammarlens}[title={What you've built}]
One more word for saying what happens if.
\end{grammarlens}

\subsection*{Your turn}
\begin{itemize}
\raggedright
  \item \emph{Say it:} \emph{nirdhāriñcu}
  \item \emph{Say it:} \emph{nirdhāriñcu}, once more
  \item \emph{Say it:} \emph{pōlīsulu ā vārtanu nirdhāriñcāru.} — The police confirmed that news.
  \item \emph{Recall:} say the Telugu for vague, then the Telugu for serious, then say \emph{nirdhāriñcu} again
\end{itemize}

\begin{tcolorbox}[breakable,colback=teal!4,colframe=teal!35!black,title={Before you move on}]
What does \te{నిర్ధారించు} mean? (\textquotedblleft{}To confirm; to ascertain\textquotedblright{}.) Say it once more.
\end{tcolorbox}
\section[(dialogue)]{Rules, risks and conditions}

\label{lesson:TE-R437-review-430-437}



\begin{quote}
Say the Telugu for to strengthen and to weaken.
\end{quote}

\subsection*{Your turn}
Say these aloud:

\begin{itemize}
\raggedright
  \item to strengthen, to weaken, balance, to control, a limit
  \item a standard, a level, a trend, effective, to test out
  \item ethics, an ideal, prestige, credibility, awareness
  \item endeavour, a challenge, a strategy, dominance, to be deprived of
  \item a system, a policy, a state of affairs, an occasion, to be formed
  \item an intellectual, a scholar, philosophy, a reformer, to inspire
  \item to take part, to contribute, to rally, to attract, to divide
  \item complicated, uncomplicated, vague, serious, to confirm
\end{itemize}

\begin{tcolorbox}[breakable,colback=teal!4,colframe=teal!35!black,title={Before you move on}]
To strengthen? (\textbf{\te{బలపరచు}}, \emph{balaparacu}.) To weaken? (\textbf{\te{బలహీనపరచు}}, \emph{balahīnaparacu}.) Balance? (\textbf{\te{సమతుల్యం}}, \emph{samatulyaṁ}.) To control? (\textbf{\te{నియంత్రించు}}, \emph{niyantriñcu}.) A limit? (\textbf{\te{పరిమితి}}, \emph{parimiti}.) A standard? (\textbf{\te{ప్రమాణం}}, \emph{pramāṇaṁ}.) A level? (\textbf{\te{స్థాయి}}, \emph{sthāyi}.) A trend? (\textbf{\te{ధోరణి}}, \emph{dhōraṇi}.) Effective? (\textbf{\te{సమర్థవంతమైన}}, \emph{samarthavantamaina}.) To test out? (\textbf{\te{పరీక్షించు}}, \emph{parīkṣiñcu}.) Ethics? (\textbf{\te{నీతి}}, \emph{nīti}.) An ideal? (\textbf{\te{ఆదర్శం}}, \emph{ādarśaṁ}.) Prestige? (\textbf{\te{ప్రతిష్ఠ}}, \emph{pratiṣṭha}.) Credibility? (\textbf{\te{విశ్వసనీయత}}, \emph{viśvasanīyata}.) Awareness? (\textbf{\te{అవగాహన}}, \emph{avagāhana}.) Endeavour? (\textbf{\te{కృషి}}, \emph{kṛṣi}.) A challenge? (\textbf{\te{సవాలు}}, \emph{savālu}.) A strategy? (\textbf{\te{వ్యూహం}}, \emph{vyūhaṁ}.) Dominance? (\textbf{\te{ఆధిపత్యం}}, \emph{ādhipatyaṁ}.) To be deprived of? (\textbf{\te{కోల్పోవు}}, \emph{kōlpōvu}.) A system? (\textbf{\te{వ్యవస్థ}}, \emph{vyavastha}.) A policy? (\textbf{\te{విధానం}}, \emph{vidhānaṁ}.) A state of affairs? (\textbf{\te{స్థితి}}, \emph{sthiti}.) An occasion? (\textbf{\te{సందర్భం}}, \emph{sandarbhaṁ}.) To be formed? (\textbf{\te{ఏర్పడు}}, \emph{ērpaḍu}.) An intellectual? (\textbf{\te{మేధావి}}, \emph{mēdhāvi}.) A scholar? (\textbf{\te{పండితుడు}}, \emph{paṇḍituḍu}.) Philosophy? (\textbf{\te{తత్వశాస్త్రం}}, \emph{tatvaśāstraṁ}.) A reformer? (\textbf{\te{సంస్కర్త}}, \emph{saṁskarta}.) To inspire? (\textbf{\te{ప్రేరేపించు}}, \emph{prērēpiñcu}.) To take part? (\textbf{\te{పాల్గొను}}, \emph{pālgonu}.) To contribute? (\textbf{\te{తోడ్పడు}}, \emph{tōḍpaḍu}.) To rally? (\textbf{\te{కూడగట్టు}}, \emph{kūḍagaṭṭu}.) To attract? (\textbf{\te{ఆకర్షించు}}, \emph{ākarṣiñcu}.) To divide? (\textbf{\te{విభజించు}}, \emph{vibhajiñcu}.) Complicated? (\textbf{\te{సంక్లిష్టమైన}}, \emph{saṅkliṣṭamaina}.) Uncomplicated? (\textbf{\te{సరళమైన}}, \emph{saraḷamaina}.) Vague? (\textbf{\te{అస్పష్టమైన}}, \emph{aspaṣṭamaina}.) Serious? (\textbf{\te{తీవ్రమైన}}, \emph{tīvramaina}.) To confirm? (\textbf{\te{నిర్ధారించు}}, \emph{nirdhāriñcu}.)
\end{tcolorbox}
\section[(dialogue)]{Hard or Plain — once more}

\label{lesson:TE-R437-review-final-437}



\begin{quote}
Say the Telugu for complicated and uncomplicated.
\end{quote}

\subsection*{Your turn}
\begin{itemize}
\raggedright
  \item \emph{Say it:} complicated, uncomplicated, vague, serious, to confirm
\end{itemize}

\begin{tcolorbox}[breakable,colback=teal!4,colframe=teal!35!black,title={Before you move on}]
Complicated? (\textbf{\te{సంక్లిష్టమైన}}, \emph{saṅkliṣṭamaina}.) Uncomplicated? (\textbf{\te{సరళమైన}}, \emph{saraḷamaina}.) Vague? (\textbf{\te{అస్పష్టమైన}}, \emph{aspaṣṭamaina}.) Serious? (\textbf{\te{తీవ్రమైన}}, \emph{tīvramaina}.) To confirm? (\textbf{\te{నిర్ధారించు}}, \emph{nirdhāriñcu}.)
\end{tcolorbox}
